\section{Companion blocks and similarity}

For a monic polynomial
\[
  p=X^d+a_{d-1}X^{d-1}+\cdots+a_0,
\]
the executable companion matrix uses the quotient power-basis convention:
ones lie immediately below the diagonal, and the last column is
$(-a_0,\ldots,-a_{d-1})^\mathsf{T}$.  We identify this matrix with
multiplication by the quotient root in mathlib's power basis and prove
\[
  \chi_{\operatorname{companion}(p)}=p.
\]
This fixes the transpose-sensitive companion convention used in the classical
rational-canonical construction \citep[Chapter VI]{macduffee1943rcf}.

Coordinatewise multiplication by $X$ on the cyclic product is block
diagonal.  Transport through the basis above proves that the transformed
matrix is exactly the dependent block family
\[
  \RCF(A)=
  \operatorname{blockDiag}
    \bigl(\operatorname{companion}(d_i)\bigr),
\]
reindexed by the sum-of-degrees equivalence.  Unit factors have degree zero
and therefore contribute empty blocks.

The public \code{RationalCanonicalResult A} stores only:
\[
  P,\qquad \text{a certificate for }P^{-1},\qquad
  P^{-1}AP=\RCF(A).
\]
The inverse certificate contains both $P^{-1}P=I$ and $PP^{-1}=I$.  Factors
and blocks are derived from $A$ rather than duplicated in the structure.

For basis independence, \code{SimilarityCertificate A B} records a matrix
$P$, its chosen inverse, and $P^{-1}AP=B$.  Applying polynomial constants
gives a two-sided equivalence
\[
  C(P^{-1})(XI-A)C(P)=XI-B.
\]
Composing this certificate with the two strong Smith results yields a
two-sided equivalence between their canonical diagonals.  Smith uniqueness
then gives exact factor equality.  Consequently
\[
  \RCF(A)=\RCF(B).
\]
Instantiating the certificate with mathlib's change-of-basis matrices proves
that both the complete factor list and canonical block matrix of an abstract
endomorphism are independent of its finite basis.
